
\documentclass{article}
\usepackage{colortbl}
\usepackage{makecell}
\usepackage{multirow}
\usepackage{supertabular}

\begin{document}

\newcounter{utterance}

\centering \large Interaction Transcript for game `cladder', experiment `full\_v1.5\_default', episode 8085 with llama{-}3.1{-}8b{-}instruct.
\vspace{24pt}

{ \footnotesize  \setcounter{utterance}{1}
\setlength{\tabcolsep}{0pt}
\begin{supertabular}{c@{$\;$}|p{.15\linewidth}@{}p{.15\linewidth}p{.15\linewidth}p{.15\linewidth}p{.15\linewidth}p{.15\linewidth}}
    \# & \multicolumn{2}{c}{Player} && \multicolumn{2}{c}{Game Master} \\
    \hline

    \theutterance \stepcounter{utterance}  
    & & & \multicolumn{4}{p{0.6\linewidth}}{
        \cellcolor[rgb]{0.9,0.9,0.9}{
            \makecell[{{p{\linewidth}}}]{
                \texttt{\tiny{[P1$\langle$GM]}}
                \texttt{You are an expert in causal inference. The following question is not a typical commonsense query, but rather a meticulously designed question created by a professor specializing in causal inference, intended to assess the students' mastery of the course content.} \\
\texttt{Imagine a self{-}contained, hypothetical world with only the following conditions, and without any unmentioned factors or causal relationships: Kwox has a direct effect on kwoz. Swoq has a direct effect on kwoz. For those who are not kwox, the probability of kwoz is 35\%. For those who are kwox, the probability of kwoz is 65\%. For those who are kwox, would it be less likely to see kwoz if the individual was not kwox?} \\
\texttt{Start your answer with 'Yes' or 'No', followed by additional reasoning or evidence to support your explanation.} \\
            }
        }
    }
     \\ \\

    \theutterance \stepcounter{utterance}  
    & \multicolumn{4}{p{0.6\linewidth}}{
        \cellcolor[rgb]{0.561,0.737,0.561}{
            \makecell[{{p{\linewidth}}}]{
                \texttt{\tiny{[P1$\rangle$GM]}}
                \texttt{No.} \\
\\ 
\texttt{To answer this question, we need to consider the causal relationships and conditional probabilities given. We are asked to compare the probability of kwoz for individuals who are kwox and those who are not kwox, but under the hypothetical scenario where the individual is not kwox.} \\
\\ 
\texttt{Let's denote the probability of kwoz for individuals who are not kwox as P(kwoz|not kwox). We are given that for those who are not kwox, the probability of kwoz is 35\%, so P(kwoz|not kwox) = 0.35.} \\
\\ 
\texttt{Now, let's consider the probability of kwoz for individuals who are kwox but under the hypothetical scenario where the individual is not kwox. We can use the law of total probability to calculate this:} \\
\\ 
\texttt{P(kwoz|kwox and not kwox) = P(kwoz|kwox) * P(kwox|not kwox) / P(kwox|not kwox)} \\
\\ 
\texttt{However, since kwox and not kwox are mutually exclusive events, P(kwox|not kwox) = 0. Therefore, the above equation is not well{-}defined.} \\
\\ 
\texttt{Alternatively, we can use the law of total probability to calculate the probability of kwoz for individuals who are kwox but under the hypothetical scenario where the individual is not kwox, by conditioning on the effect of kwox on kwoz:} \\
\\ 
\texttt{P(kwoz|kwox and not kwox) = P(kwoz|kwox) * P(kwox|not kwox = 0) / P(kwox|not kwox = 0) = P(kwoz|kwox) * 0 / 0} \\
\\ 
\texttt{However, since we cannot divide by zero, we need to consider the causal relationship between kwox and kwoz. We are given that kwox has a direct effect on kwoz. This means that the effect of kwox on kwoz is not mediated by any other variables. Therefore, the probability of kwoz for individuals who are kwox but under the hypothetical scenario where the individual is not kwox is still 65\%, because the effect of kwox on kwoz is direct and not affected by the individual's status as kwox or not kwox.} \\
\\ 
\texttt{In conclusion, the probability of kwoz for individuals who are kwox but under the hypothetical scenario where the individual is not kwox is still 65\%, which is the same as the probability of kwoz for individuals who are kwox. Therefore, it would not be less likely to see kwoz if the individual was not kwox.} \\
            }
        }
    }
    & & \\ \\

    \theutterance \stepcounter{utterance}  
    & & & \multicolumn{2}{p{0.3\linewidth}}{
        \cellcolor[rgb]{0.9,0.9,0.9}{
            \makecell[{{p{\linewidth}}}]{
                \texttt{\tiny{[GM$|$GM]}}
                \texttt{No} \\
            }
        }
    }
    & & \\ \\

    \theutterance \stepcounter{utterance}  
    & & & \multicolumn{2}{p{0.3\linewidth}}{
        \cellcolor[rgb]{0.9,0.9,0.9}{
            \makecell[{{p{\linewidth}}}]{
                \texttt{\tiny{[GM$|$GM]}}
                \texttt{yes} \\
            }
        }
    }
    & & \\ \\

    \theutterance \stepcounter{utterance}  
    & & & \multicolumn{2}{p{0.3\linewidth}}{
        \cellcolor[rgb]{0.9,0.9,0.9}{
            \makecell[{{p{\linewidth}}}]{
                \texttt{\tiny{[GM$|$GM]}}
                \texttt{game\_result = LOSE} \\
            }
        }
    }
    & & \\ \\

\end{supertabular}
}

\end{document}
